%% Everything between the BODY markers is journal-independent: util/new_paper.py --switch
%% carries it to another journal class unchanged. Keep class-specific markup out of it.

\section{Introduction}
\label{sec:intro}

\todo{Open with the problem and why it matters, in terms a reader outside the subfield can follow. Then say what is known, citing the work that established it \citep{Hogg2010, ForemanMackey2013}, and what is not. The last paragraph says what this paper does and what each section contains, in full sentences.}

\section{Methods}
\label{sec:methods}

\todo{Describe what we did in enough detail that a competent reader could repeat it. Every number quoted here should be traceable to a script in this repository or to a note in research/.}

Number an equation only if the text refers to it, and punctuate a displayed equation as part of its sentence. The likelihood of $N$ independent measurements $y_i$ with Gaussian uncertainties $\sigma_i$ under a model $f(x;\theta)$ is
\begin{equation}
\ln \mathcal{L}(\theta) = -\frac{1}{2} \sum_{i=1}^{N} \left[ \frac{\left(y_i - f(x_i;\theta)\right)^2}{\sigma_i^2} + \ln\left(2\pi\sigma_i^2\right) \right],
\label{eq:likelihood}
\end{equation}
and we sample the posterior with an affine-invariant ensemble sampler \citep{Goodman2010, ForemanMackey2013}. Equation~(\ref{eq:likelihood}) assumes that the uncertainties are known; Section~\ref{sec:discussion} discusses what happens when they are not.

\begin{proposition}
\label{prop:example}
If every $\sigma_i = \sigma$ and $f$ is linear in $\theta$, the maximum of $\mathcal{L}$ is the ordinary least-squares solution.
\end{proposition}

\begin{proof}
With equal $\sigma_i$ the second term in Equation~(\ref{eq:likelihood}) does not depend on $\theta$, so maximizing $\ln\mathcal{L}$ minimizes $\sum_i (y_i - f(x_i;\theta))^2$.
\end{proof}

\section{Results}
\label{sec:results}

\todo{State each result once, with its uncertainty, and point to the figure or table that shows it.} Figure~\ref{fig:example} shows simulated measurements around a straight line, and Table~\ref{tab:example} lists the fitted values. % drawn by paper/figures/make_example.py

\begin{figure}
\centering
\includegraphics[width=\columnwidth]{figures/example.pdf}
\caption{A caption says what is plotted, what the symbols and colors mean, and what the reader should take away. Points: simulated measurements with $1\sigma$ uncertainties. Line: the input model.}
\label{fig:example}
\end{figure}

\begin{table}
\centering
\caption{Example table. Units go in the column heads, not in the cells.}
\label{tab:example}
\begin{tabular}{lcc}
\toprule
Quantity & Value & Uncertainty \\
\midrule
Slope $m$ & 0.52 & 0.03 \\
Intercept $b$ (dex) & $-1.10$ & 0.07 \\
\bottomrule
\end{tabular}
\end{table}

\section{Discussion}
\label{sec:discussion}

\todo{What the results mean, what they do not show, and how they compare with earlier work.}

\section{Conclusions}
\label{sec:conclusions}

\todo{The findings, in the order of their importance. No new material.}

\begin{acknowledgments}
\todo{Funding, data providers, and the people who helped.}
\end{acknowledgments}

\software{NumPy \citep{Harris2020}, Matplotlib \citep{Hunter2007}, Astropy \citep{Astropy2022}}

\appendix

\section{Derivations}
\label{app:derivations}

\todo{Material that a careful reader needs but that interrupts the argument. Delete this appendix if it is empty.}
